\chapter{Limitations and Future Work}
\label{ch:limitations}

\section{Limitations}
The thesis states its limits plainly, since doing so is part of an honest benchmark.

\begin{itemize}[leftmargin=1.4em]
  \item \textbf{Single grader, no inter-annotator agreement.} A second grader could not be obtained
  within the scope of this work, so there is no measured human agreement ceiling. All results are
  agreement \emph{with this teacher}, and may inherit that grader's biases. Bootstrap confidence
  intervals partially mitigate, but do not remove, this limitation.
  \item \textbf{Question leakage in the headline split.} Because only 41 distinct questions exist, the
  random split shares questions between train and test. The leakage is now \emph{measured}
  (\S\ref{sec:leakage}) rather than merely flagged, and both seen- and unseen-question numbers are
  reported; but the seen-question headline numbers should be read with this in mind.
  \item \textbf{Small corpus.} With roughly 895--1{,}184 answers and few questions, the
  unseen-question test sets are tiny ($\sim$6 questions) and high-variance, and the corpus covers four
  subjects from two schools rather than a national sample.
  \item \textbf{Partial cross-pillar comparability.} The pillars were evaluated on different dataset
  variants (895/909/1184) and only the classical and BiLSTM models share a test set, so cross-pillar
  QWK differences should be read as indicative rather than head-to-head.
  \item \textbf{Explainability scope.} The encoder and LLM LOO faithfulness rows require GPU hardware
  and are \emph{HPC-pending}. The plausibility measure is an automatic reference-overlap proxy, not a
  human-rationale study.
  \item \textbf{No formal field study.} The live prototype demonstrates feasibility, but its
  classroom value has not yet been evaluated with teachers in a controlled study.
\end{itemize}

\section{Future Work}
Several directions follow naturally and would address the limitations above.

\begin{itemize}[leftmargin=1.4em]
  \item \textbf{A second grader and a question-held-out benchmark.} Adding a second annotator would
  give a true agreement ceiling, and re-running the \emph{full} grid (all four pillars and the LLM)
  under the question-held-out split would turn the leakage analysis into the primary, honest
  generalisation benchmark.
  \item \textbf{Khmer-specific pretraining.} Continued pretraining of a multilingual encoder on Khmer
  text, or a Khmer-native encoder, may close the gap left by the limited Khmer seen during generic
  multilingual pretraining.
  \item \textbf{Full-pillar LOO faithfulness.} Running the encoder and LLM LOO faithfulness rows on
  GPU hardware would extend the verified faithful explanation to all four pillars.
  \item \textbf{Cross-lingual transfer.} Transferring scoring ability from high-resource languages to
  Khmer (for example, via multilingual representations or cross-lingual contrastive
  learning~\citep{li2024crosslingual}) could help in the genuinely low-data, unseen-question regime.
  \item \textbf{Feedback quality.} The prototype's written feedback should be evaluated for
  pedagogical usefulness in classroom settings, building on rubric-aligned feedback
  work~\citep{condor2024,nam2024grading}.
  \item \textbf{A teacher field study.} A controlled study of how Cambodian teachers use, trust, and
  benefit from the prototype would test the explainability contribution where it ultimately matters.
\end{itemize}
